\documentclass[UTF8,12pt]{ctexart}
\usepackage[a4paper,margin=24mm]{geometry}
\usepackage{amsmath,amssymb,bm,graphicx,adjustbox}
\newcommand{\fitmath}[1]{\begin{adjustbox}{max width=\linewidth}$\displaystyle #1$\end{adjustbox}}
\setlength{\parindent}{0pt}
\setlength{\parskip}{.7em}
\sloppy
\allowdisplaybreaks
\begin{document}
\section*{2006 年考研数学三 · 参考答案}
\subsection*{2006年真题参考答案}

\subsection*{一、填空题}

（1）\(\displaystyle 1\)。（2）\(\displaystyle 2e^3\)。（3）\(\displaystyle 4\,\mathrm dx-2\,\mathrm dy\)。（4）\(\displaystyle 2\)。（5）\(\displaystyle \dfrac{\displaystyle 1}{\displaystyle 9}\)。（6）\(\displaystyle 2\)。


\subsection*{二、选择题}

（7）A。（8）C。（9）D。（10）B。（11）D。（12）A。（13）B。（14）A。


\subsection*{三、解答题}

（15）（I）\(\displaystyle g(x)=\dfrac{\displaystyle 1}{\displaystyle x}-\dfrac{\displaystyle 1-\pi x}{\displaystyle \arctan x}\)。（II）\(\displaystyle \pi\)。（16）\(\displaystyle \dfrac{\displaystyle 2}{\displaystyle 9}\)。


（17）证明略。（可考虑 \(\displaystyle f(x)=x\sin x+2\cos x+\pi x\)，\(\displaystyle x\in[0,\allowbreak \pi]\)，利用导数分析 \(\displaystyle f(x)\) 的单调性。）


（18）（I）\(\displaystyle y=ax^2-ax\)。（II）\(\displaystyle a=2\)。


（19）收敛域为 \(\displaystyle [-1,\allowbreak 1]\)；\(\displaystyle S(x)=2x^2\arctan x-x\ln(1+x^2)\)，\(\displaystyle x\in[-1,\allowbreak 1]\)。


（20）当 \(\displaystyle a=0\) 时，\(\displaystyle \alpha_1,\allowbreak \alpha_2,\allowbreak \alpha_3,\allowbreak \alpha_4\) 线性相关，\(\displaystyle \alpha_1\) 为一个极大线性无关组，且 \(\displaystyle \alpha_2=2\alpha_1\)，\(\displaystyle \alpha_3=3\alpha_1\)，\(\displaystyle \alpha_4=4\alpha_1\)。当 \(\displaystyle a=-10\) 时，向量组线性相关，\(\displaystyle \alpha_2,\allowbreak \alpha_3,\allowbreak \alpha_4\) 为一个极大线性无关组，且 \(\displaystyle \alpha_1=-\alpha_2-\alpha_3-\alpha_4\)。


（21）（I）\(\displaystyle A\) 的特征值为 \(\displaystyle 0,\allowbreak 0,\allowbreak 3\)。属于特征值0的特征向量为 \(\displaystyle k_1\alpha_1+k_2\alpha_2\)，其中 \(\displaystyle k_1,\allowbreak k_2\) 为不全为零的任意常数；属于特征值3的特征向量为 \(\displaystyle k(1,\allowbreak 1,\allowbreak 1)^{\mathrm T}\)，其中 \(\displaystyle k\) 为任意非零常数。


（21）（II）\(\displaystyle Q=\begin{pmatrix}-\dfrac{\displaystyle 1}{\displaystyle \sqrt6}&-\dfrac{\displaystyle 1}{\displaystyle \sqrt2}&\dfrac{\displaystyle 1}{\displaystyle \sqrt3}\\\dfrac{\displaystyle 2}{\displaystyle \sqrt6}&0&\dfrac{\displaystyle 1}{\displaystyle \sqrt3}\\-\dfrac{\displaystyle 1}{\displaystyle \sqrt6}&\dfrac{\displaystyle 1}{\displaystyle \sqrt2}&\dfrac{\displaystyle 1}{\displaystyle \sqrt3}\end{pmatrix}\)，\(\displaystyle \Lambda=\begin{pmatrix}0&0&0\\0&0&0\\0&0&3\end{pmatrix}\)。（III）\(\displaystyle A=\begin{pmatrix}1&1&1\\1&1&1\\1&1&1\end{pmatrix}\)，\(\displaystyle (A-\dfrac{\displaystyle 3}{\displaystyle 2}E)^6=(\dfrac{\displaystyle 3}{\displaystyle 2})^6E\)。


（22）（I）\(\displaystyle f_Y(y)=\begin{cases}\dfrac{\displaystyle 3}{\displaystyle 8\sqrt y},\allowbreak &0<y<1,\allowbreak \\\dfrac{\displaystyle 1}{\displaystyle 8\sqrt y},\allowbreak &1\le y<4,\allowbreak \\0,\allowbreak &\text{其他}.\end{cases}\)（II）\(\displaystyle \operatorname{Cov}(X,\allowbreak Y)=\dfrac{\displaystyle 2}{\displaystyle 3}\)。（III）\(\displaystyle \dfrac{\displaystyle 1}{\displaystyle 4}\)。


（23）（I）\(\displaystyle \hat\theta=\dfrac{\displaystyle 3}{\displaystyle 2}-\bar X\)。（II）\(\displaystyle \hat\theta=\dfrac{\displaystyle N}{\displaystyle n}\)。

\end{document}
